
\documentclass{article}
\usepackage{colortbl}
\usepackage{makecell}
\usepackage{multirow}
\usepackage{supertabular}

\begin{document}

\newcounter{utterance}

\centering \large Interaction Transcript for game `dond', experiment `semi\_en', episode 12 with Qwen3{-}32B{-}t0.0.
\vspace{24pt}

{ \footnotesize  \setcounter{utterance}{1}
\setlength{\tabcolsep}{0pt}
\begin{supertabular}{c@{$\;$}|p{.15\linewidth}@{}p{.15\linewidth}p{.15\linewidth}p{.15\linewidth}p{.15\linewidth}p{.15\linewidth}}
    \# & $\;$A & \multicolumn{4}{c}{Game Master} & $\;\:$B\\
    \hline

    \theutterance \stepcounter{utterance}  
    & & \multicolumn{4}{p{0.6\linewidth}}{
        \cellcolor[rgb]{0.9,0.9,0.9}{
            \makecell[{{p{\linewidth}}}]{
                \texttt{\tiny{[P1$\langle$GM]}}
                \texttt{You are playing a negotiation game in which you have to agree on how to divide a set of items among you and another player.} \\
\\ 
\texttt{Rules:} \\
\texttt{(a) You and the other player are given a set of items. Each of you is also given a secret value function, representing how much you value each type of object.} \\
\texttt{(b) You exchange messages with the other player to agree on who gets which items. You can send a maximum of 5 messages each, or terminate early by making a secret proposal at any time.} \\
\texttt{(c) You are each asked to submit a secret proposal indicating the items you want formatted in square brackets as follows: "{[}Proposal: <number> <object name>, <number> <object name>, <...>{]}"} \\
\texttt{(d) If your proposals are complementary, i.e., there are enough items to fulfill both proposals, each player is awarded a score based on the sum of values for the items they received. Otherwise, both of you get zero points.} \\
\texttt{(e) Your goal is to maximize the score you receive.} \\
\\ 
\texttt{Let us start.} \\
\\ 
\texttt{The set of available items is:} \\
\\ 
\texttt{1 lamp, 3 thimbles, 1 rope} \\
\\ 
\texttt{Your secret value function is:} \\
\\ 
\texttt{lamp: 3, thimble: 2, rope: 1} \\
\\ 
\texttt{IMPORTANT: Your messages, unless it is the secret proposal, are directly transmitted to the other player, so do not include any response to the rules or text announcing your message. To make a secret proposal, use the indicated format. Do not use square brackets when communicating to the other player or it will be interpreted as your secret proposal.} \\
            }
        }
    }
    & \\ \\

    \theutterance \stepcounter{utterance}  
    & \multicolumn{4}{p{0.6\linewidth}}{
        \cellcolor[rgb]{0.9,0.9,0.9}{
            \makecell[{{p{\linewidth}}}]{
                \texttt{\tiny{[P1$\rangle$GM]}}
                \texttt{<think>} \\
\texttt{Okay, let's see. I need to negotiate with another player to divide the items: 1 lamp, 3 thimbles, and 1 rope. My value function is lamp: 3, thimble: 2, rope: 1. The other player's value function is secret, so I don't know what they value. The goal is to maximize my own score, but we both have to have complementary proposals. If our proposals don't conflict, we each get our items. If they do, we both get zero.} \\
\\ 
\texttt{First, I should figure out what the other player might want. Since the items are limited, maybe they value the lamp or thimbles highly. But without knowing their values, I have to make assumptions. Let's think about possible scenarios.} \\
\\ 
\texttt{If I take the lamp, that's 3 points for me. The other player might want the lamp too. But if they don't, maybe they want thimbles. Since there are 3 thimbles, maybe I can take some and leave some for them. Alternatively, if I take the rope, it's only 1 point for me, but maybe the other player values it more. But I don't know.} \\
\\ 
\texttt{Since my highest value is the lamp, I should probably try to get the lamp. But I need to make sure that the other player can get their desired items without overlapping. Let's consider possible proposals.} \\
\\ 
\texttt{Option 1: I take the lamp. Then the remaining items are 3 thimbles and 1 rope. If the other player wants thimbles, they could take all 3, but maybe they need some other combination. Alternatively, if they want the rope, they can take it. But since I don't know their values, I need to make a proposal that allows them to have something without conflict.} \\
\\ 
\texttt{Option 2: I take two thimbles (2*2=4 points) and the rope (1 point), totaling 5. But then the other player would have the lamp and one thimble. But if they value the lamp highly, maybe they would prefer that. But again, I don't know.} \\
\\ 
\texttt{Wait, but the problem is that if I take the lamp, the other player can't take it. So if they also want the lamp, our proposals would conflict. But since I don't know their values, I have to make a proposal that's likely to be complementary.} \\
\\ 
\texttt{Another} \\
            }
        }
    }
    & & \\ \\

    \theutterance \stepcounter{utterance}  
    & & & \multicolumn{2}{p{0.3\linewidth}}{
        \cellcolor[rgb]{0.9,0.9,0.9}{
            \makecell[{{p{\linewidth}}}]{
                \texttt{\tiny{[GM$|$GM]}}
                \texttt{continue} \\
            }
        }
    }
    & & \\ \\

    \theutterance \stepcounter{utterance}  
    & & & \multicolumn{2}{p{0.3\linewidth}}{
        \cellcolor[rgb]{0.9,0.9,0.9}{
            \makecell[{{p{\linewidth}}}]{
                \texttt{\tiny{[GM$|$GM]}}
                \texttt{<think>} \\
\texttt{Okay, let's see. I need to negotiate with another player to divide the items: 1 lamp, 3 thimbles, and 1 rope. My value function is lamp: 3, thimble: 2, rope: 1. The other player's value function is secret, so I don't know what they value. The goal is to maximize my own score, but we both have to have complementary proposals. If our proposals don't conflict, we each get our items. If they do, we both get zero.} \\
\\ 
\texttt{First, I should figure out what the other player might want. Since the items are limited, maybe they value the lamp or thimbles highly. But without knowing their values, I have to make assumptions. Let's think about possible scenarios.} \\
\\ 
\texttt{If I take the lamp, that's 3 points for me. The other player might want the lamp too. But if they don't, maybe they want thimbles. Since there are 3 thimbles, maybe I can take some and leave some for them. Alternatively, if I take the rope, it's only 1 point for me, but maybe the other player values it more. But I don't know.} \\
\\ 
\texttt{Since my highest value is the lamp, I should probably try to get the lamp. But I need to make sure that the other player can get their desired items without overlapping. Let's consider possible proposals.} \\
\\ 
\texttt{Option 1: I take the lamp. Then the remaining items are 3 thimbles and 1 rope. If the other player wants thimbles, they could take all 3, but maybe they need some other combination. Alternatively, if they want the rope, they can take it. But since I don't know their values, I need to make a proposal that allows them to have something without conflict.} \\
\\ 
\texttt{Option 2: I take two thimbles (2*2=4 points) and the rope (1 point), totaling 5. But then the other player would have the lamp and one thimble. But if they value the lamp highly, maybe they would prefer that. But again, I don't know.} \\
\\ 
\texttt{Wait, but the problem is that if I take the lamp, the other player can't take it. So if they also want the lamp, our proposals would conflict. But since I don't know their values, I have to make a proposal that's likely to be complementary.} \\
\\ 
\texttt{Another} \\
            }
        }
    }
    & & \\ \\

    \theutterance \stepcounter{utterance}  
    & & \multicolumn{4}{p{0.6\linewidth}}{
        \cellcolor[rgb]{0.9,0.9,0.9}{
            \makecell[{{p{\linewidth}}}]{
                \texttt{\tiny{[GM$\rangle$P2]}}
                \texttt{You are playing a negotiation game in which you have to agree on how to divide a set of items among you and another player.} \\
\\ 
\texttt{Rules:} \\
\texttt{(a) You and the other player are given a set of items. Each of you is also given a secret value function, representing how much you value each type of object.} \\
\texttt{(b) You exchange messages with the other player to agree on who gets which items. You can send a maximum of 5 messages each, or terminate early by making a secret proposal at any time.} \\
\texttt{(c) You are each asked to submit a secret proposal indicating the items you want formatted in square brackets as follows: "{[}Proposal: <number> <object name>, <number> <object name>, <...>{]}"} \\
\texttt{(d) If your proposals are complementary, i.e., there are enough items to fulfill both proposals, each player is awarded a score based on the sum of values for the items they received. Otherwise, both of you get zero points.} \\
\texttt{(e) Your goal is to maximize the score you receive.} \\
\\ 
\texttt{Let us start.} \\
\\ 
\texttt{The set of available items is:} \\
\\ 
\texttt{1 lamp, 3 thimbles, 1 rope} \\
\\ 
\texttt{Your secret value function is:} \\
\\ 
\texttt{lamp: 3, thimble: 1, rope: 4} \\
\\ 
\texttt{IMPORTANT: Your messages, unless it is the secret proposal, are directly transmitted to the other player, so do not include any response to the rules or text announcing your message. To make a secret proposal, use the indicated format. Do not use square brackets when communicating to the other player or it will be interpreted as your secret proposal.} \\
\\ 
\\ 
\texttt{<think>} \\
\texttt{Okay, let's see. I need to negotiate with another player to divide the items: 1 lamp, 3 thimbles, and 1 rope. My value function is lamp: 3, thimble: 2, rope: 1. The other player's value function is secret, so I don't know what they value. The goal is to maximize my own score, but we both have to have complementary proposals. If our proposals don't conflict, we each get our items. If they do, we both get zero.} \\
\\ 
\texttt{First, I should figure out what the other player might want. Since the items are limited, maybe they value the lamp or thimbles highly. But without knowing their values, I have to make assumptions. Let's think about possible scenarios.} \\
\\ 
\texttt{If I take the lamp, that's 3 points for me. The other player might want the lamp too. But if they don't, maybe they want thimbles. Since there are 3 thimbles, maybe I can take some and leave some for them. Alternatively, if I take the rope, it's only 1 point for me, but maybe the other player values it more. But I don't know.} \\
\\ 
\texttt{Since my highest value is the lamp, I should probably try to get the lamp. But I need to make sure that the other player can get their desired items without overlapping. Let's consider possible proposals.} \\
\\ 
\texttt{Option 1: I take the lamp. Then the remaining items are 3 thimbles and 1 rope. If the other player wants thimbles, they could take all 3, but maybe they need some other combination. Alternatively, if they want the rope, they can take it. But since I don't know their values, I need to make a proposal that allows them to have something without conflict.} \\
\\ 
\texttt{Option 2: I take two thimbles (2*2=4 points) and the rope (1 point), totaling 5. But then the other player would have the lamp and one thimble. But if they value the lamp highly, maybe they would prefer that. But again, I don't know.} \\
\\ 
\texttt{Wait, but the problem is that if I take the lamp, the other player can't take it. So if they also want the lamp, our proposals would conflict. But since I don't know their values, I have to make a proposal that's likely to be complementary.} \\
\\ 
\texttt{Another} \\
            }
        }
    }
    & \\ \\

    \theutterance \stepcounter{utterance}  
    & & & \multicolumn{4}{p{0.6\linewidth}}{
        \cellcolor[rgb]{0.9,0.9,0.9}{
            \makecell[{{p{\linewidth}}}]{
                \texttt{\tiny{[GM$\langle$P2]}}
                \texttt{<think>} \\
\texttt{Okay, let's try to figure out the best way to approach this. The available items are 1 lamp, 3 thimbles, and 1 rope. My values are lamp: 3, thimble: 2, rope: 1. The other player's values are secret, so I need to make a proposal that maximizes my points while allowing them to have a complementary set.} \\
\\ 
\texttt{First, I should consider what the other player might value. Since there's only one lamp and one rope, they might be high{-}value items for them too. But with three thimbles, maybe they can take some without conflicting. Let's think about possible scenarios.} \\
\\ 
\texttt{If I take the lamp (3 points), then the remaining items are 3 thimbles and 1 rope. The other player could take all 3 thimbles and the rope. If their value for thimbles is low, maybe they'd prefer the rope. But if they value thimbles more, taking all three might be good. However, since I don't know their values, I need to make a proposal that's flexible.} \\
\\ 
\texttt{Alternatively, if I take two thimbles (2*2=4) and the rope (1), that's 5 points. Then the other player could take the lamp and one thimble. But if they value the lamp highly, that might work. But if they also want thimbles, maybe they'd prefer more of them. However, there are only three, so if I take two, they can take one. But maybe they want the lamp and the rope instead.} \\
\\ 
\texttt{Another option is to take the lamp and one thimble (3+2=5). Then the other player could take two thimbles and the rope. But again, depends on their values.} \\
\\ 
\texttt{Wait, the key is to make a proposal that doesn't conflict with possible other proposals. Since the other player's values are unknown, I need to choose a proposal that allows them to have a set that doesn't overlap with mine. The safest way might be to take the lamp and leave the thimbles and rope for them. Because if they want the lamp, we'd conflict, but if they don't, they can take other items. Alternatively, if I take the rope, which I value low, maybe that leaves the lamp for them if they value it.} \\
\\ 
\texttt{But my highest value is the lamp. So maybe I should go} \\
            }
        }
    }
     \\ \\

    \theutterance \stepcounter{utterance}  
    & & & \multicolumn{2}{p{0.3\linewidth}}{
        \cellcolor[rgb]{0.9,0.9,0.9}{
            \makecell[{{p{\linewidth}}}]{
                \texttt{\tiny{[GM$|$GM]}}
                \texttt{continue} \\
            }
        }
    }
    & & \\ \\

    \theutterance \stepcounter{utterance}  
    & & & \multicolumn{2}{p{0.3\linewidth}}{
        \cellcolor[rgb]{0.9,0.9,0.9}{
            \makecell[{{p{\linewidth}}}]{
                \texttt{\tiny{[GM$|$GM]}}
                \texttt{<think>} \\
\texttt{Okay, let's try to figure out the best way to approach this. The available items are 1 lamp, 3 thimbles, and 1 rope. My values are lamp: 3, thimble: 2, rope: 1. The other player's values are secret, so I need to make a proposal that maximizes my points while allowing them to have a complementary set.} \\
\\ 
\texttt{First, I should consider what the other player might value. Since there's only one lamp and one rope, they might be high{-}value items for them too. But with three thimbles, maybe they can take some without conflicting. Let's think about possible scenarios.} \\
\\ 
\texttt{If I take the lamp (3 points), then the remaining items are 3 thimbles and 1 rope. The other player could take all 3 thimbles and the rope. If their value for thimbles is low, maybe they'd prefer the rope. But if they value thimbles more, taking all three might be good. However, since I don't know their values, I need to make a proposal that's flexible.} \\
\\ 
\texttt{Alternatively, if I take two thimbles (2*2=4) and the rope (1), that's 5 points. Then the other player could take the lamp and one thimble. But if they value the lamp highly, that might work. But if they also want thimbles, maybe they'd prefer more of them. However, there are only three, so if I take two, they can take one. But maybe they want the lamp and the rope instead.} \\
\\ 
\texttt{Another option is to take the lamp and one thimble (3+2=5). Then the other player could take two thimbles and the rope. But again, depends on their values.} \\
\\ 
\texttt{Wait, the key is to make a proposal that doesn't conflict with possible other proposals. Since the other player's values are unknown, I need to choose a proposal that allows them to have a set that doesn't overlap with mine. The safest way might be to take the lamp and leave the thimbles and rope for them. Because if they want the lamp, we'd conflict, but if they don't, they can take other items. Alternatively, if I take the rope, which I value low, maybe that leaves the lamp for them if they value it.} \\
\\ 
\texttt{But my highest value is the lamp. So maybe I should go} \\
            }
        }
    }
    & & \\ \\

    \theutterance \stepcounter{utterance}  
    & & \multicolumn{4}{p{0.6\linewidth}}{
        \cellcolor[rgb]{0.9,0.9,0.9}{
            \makecell[{{p{\linewidth}}}]{
                \texttt{\tiny{[P1$\langle$GM]}}
                \texttt{<think>} \\
\texttt{Okay, let's try to figure out the best way to approach this. The available items are 1 lamp, 3 thimbles, and 1 rope. My values are lamp: 3, thimble: 2, rope: 1. The other player's values are secret, so I need to make a proposal that maximizes my points while allowing them to have a complementary set.} \\
\\ 
\texttt{First, I should consider what the other player might value. Since there's only one lamp and one rope, they might be high{-}value items for them too. But with three thimbles, maybe they can take some without conflicting. Let's think about possible scenarios.} \\
\\ 
\texttt{If I take the lamp (3 points), then the remaining items are 3 thimbles and 1 rope. The other player could take all 3 thimbles and the rope. If their value for thimbles is low, maybe they'd prefer the rope. But if they value thimbles more, taking all three might be good. However, since I don't know their values, I need to make a proposal that's flexible.} \\
\\ 
\texttt{Alternatively, if I take two thimbles (2*2=4) and the rope (1), that's 5 points. Then the other player could take the lamp and one thimble. But if they value the lamp highly, that might work. But if they also want thimbles, maybe they'd prefer more of them. However, there are only three, so if I take two, they can take one. But maybe they want the lamp and the rope instead.} \\
\\ 
\texttt{Another option is to take the lamp and one thimble (3+2=5). Then the other player could take two thimbles and the rope. But again, depends on their values.} \\
\\ 
\texttt{Wait, the key is to make a proposal that doesn't conflict with possible other proposals. Since the other player's values are unknown, I need to choose a proposal that allows them to have a set that doesn't overlap with mine. The safest way might be to take the lamp and leave the thimbles and rope for them. Because if they want the lamp, we'd conflict, but if they don't, they can take other items. Alternatively, if I take the rope, which I value low, maybe that leaves the lamp for them if they value it.} \\
\\ 
\texttt{But my highest value is the lamp. So maybe I should go} \\
            }
        }
    }
    & \\ \\

    \theutterance \stepcounter{utterance}  
    & \multicolumn{4}{p{0.6\linewidth}}{
        \cellcolor[rgb]{0.9,0.9,0.9}{
            \makecell[{{p{\linewidth}}}]{
                \texttt{\tiny{[P1$\rangle$GM]}}
                \texttt{<think>} \\
\texttt{Okay, let's try to figure out the best way to approach this. The available items are 1 lamp, 3 thimbles, and 1 rope. My values are lamp: 3, thimble: 2, rope: 1. The other player's values are secret, so I need to make a proposal that maximizes my points while allowing them to have a complementary set.} \\
\\ 
\texttt{First, I should consider what the other player might value. Since there's only one lamp and one rope, they might be high{-}value items for them too. But with three thimbles, maybe they can take some without conflicting. Let's think about possible scenarios.} \\
\\ 
\texttt{If I take the lamp (3 points), then the remaining items are 3 thimbles and 1 rope. The other player could take all 3 thimbles and the rope. If their value for thimbles is low, maybe they'd prefer the rope. But if they value thimbles more, taking all three might be good. However, since I don't know their values, I need to make a proposal that's flexible.} \\
\\ 
\texttt{Alternatively, if I take two thimbles (2*2=4) and the rope (1), that's 5 points. Then the other player could take the lamp and one thimble. But if they value the lamp highly, that might work. But if they also want thimbles, maybe they'd prefer more of them. However, there are only three, so if I take two, they can take one. But maybe they want the lamp and the rope instead.} \\
\\ 
\texttt{Another option is to take the lamp and one thimble (3+2=5). Then the other player could take two thimbles and the rope. But again, depends on their values.} \\
\\ 
\texttt{Wait, the key is to make a proposal that doesn't conflict with possible other proposals. Since the other player's values are unknown, I need to choose a proposal that allows them to have a set that doesn't overlap with mine. The safest way might be to take the lamp and leave the thimbles and rope for them. Because if they want the lamp, we'd conflict, but if they don't, they can take other items. Alternatively, if I take the rope, which I value low, maybe that leaves the lamp for them if they value it.} \\
\\ 
\texttt{But my highest value is the lamp. So maybe I should go} \\
            }
        }
    }
    & & \\ \\

    \theutterance \stepcounter{utterance}  
    & & & \multicolumn{2}{p{0.3\linewidth}}{
        \cellcolor[rgb]{0.9,0.9,0.9}{
            \makecell[{{p{\linewidth}}}]{
                \texttt{\tiny{[GM$|$GM]}}
                \texttt{continue} \\
            }
        }
    }
    & & \\ \\

    \theutterance \stepcounter{utterance}  
    & & & \multicolumn{2}{p{0.3\linewidth}}{
        \cellcolor[rgb]{0.9,0.9,0.9}{
            \makecell[{{p{\linewidth}}}]{
                \texttt{\tiny{[GM$|$GM]}}
                \texttt{<think>} \\
\texttt{Okay, let's try to figure out the best way to approach this. The available items are 1 lamp, 3 thimbles, and 1 rope. My values are lamp: 3, thimble: 2, rope: 1. The other player's values are secret, so I need to make a proposal that maximizes my points while allowing them to have a complementary set.} \\
\\ 
\texttt{First, I should consider what the other player might value. Since there's only one lamp and one rope, they might be high{-}value items for them too. But with three thimbles, maybe they can take some without conflicting. Let's think about possible scenarios.} \\
\\ 
\texttt{If I take the lamp (3 points), then the remaining items are 3 thimbles and 1 rope. The other player could take all 3 thimbles and the rope. If their value for thimbles is low, maybe they'd prefer the rope. But if they value thimbles more, taking all three might be good. However, since I don't know their values, I need to make a proposal that's flexible.} \\
\\ 
\texttt{Alternatively, if I take two thimbles (2*2=4) and the rope (1), that's 5 points. Then the other player could take the lamp and one thimble. But if they value the lamp highly, that might work. But if they also want thimbles, maybe they'd prefer more of them. However, there are only three, so if I take two, they can take one. But maybe they want the lamp and the rope instead.} \\
\\ 
\texttt{Another option is to take the lamp and one thimble (3+2=5). Then the other player could take two thimbles and the rope. But again, depends on their values.} \\
\\ 
\texttt{Wait, the key is to make a proposal that doesn't conflict with possible other proposals. Since the other player's values are unknown, I need to choose a proposal that allows them to have a set that doesn't overlap with mine. The safest way might be to take the lamp and leave the thimbles and rope for them. Because if they want the lamp, we'd conflict, but if they don't, they can take other items. Alternatively, if I take the rope, which I value low, maybe that leaves the lamp for them if they value it.} \\
\\ 
\texttt{But my highest value is the lamp. So maybe I should go} \\
            }
        }
    }
    & & \\ \\

    \theutterance \stepcounter{utterance}  
    & & \multicolumn{4}{p{0.6\linewidth}}{
        \cellcolor[rgb]{0.9,0.9,0.9}{
            \makecell[{{p{\linewidth}}}]{
                \texttt{\tiny{[GM$\rangle$P2]}}
                \texttt{<think>} \\
\texttt{Okay, let's try to figure out the best way to approach this. The available items are 1 lamp, 3 thimbles, and 1 rope. My values are lamp: 3, thimble: 2, rope: 1. The other player's values are secret, so I need to make a proposal that maximizes my points while allowing them to have a complementary set.} \\
\\ 
\texttt{First, I should consider what the other player might value. Since there's only one lamp and one rope, they might be high{-}value items for them too. But with three thimbles, maybe they can take some without conflicting. Let's think about possible scenarios.} \\
\\ 
\texttt{If I take the lamp (3 points), then the remaining items are 3 thimbles and 1 rope. The other player could take all 3 thimbles and the rope. If their value for thimbles is low, maybe they'd prefer the rope. But if they value thimbles more, taking all three might be good. However, since I don't know their values, I need to make a proposal that's flexible.} \\
\\ 
\texttt{Alternatively, if I take two thimbles (2*2=4) and the rope (1), that's 5 points. Then the other player could take the lamp and one thimble. But if they value the lamp highly, that might work. But if they also want thimbles, maybe they'd prefer more of them. However, there are only three, so if I take two, they can take one. But maybe they want the lamp and the rope instead.} \\
\\ 
\texttt{Another option is to take the lamp and one thimble (3+2=5). Then the other player could take two thimbles and the rope. But again, depends on their values.} \\
\\ 
\texttt{Wait, the key is to make a proposal that doesn't conflict with possible other proposals. Since the other player's values are unknown, I need to choose a proposal that allows them to have a set that doesn't overlap with mine. The safest way might be to take the lamp and leave the thimbles and rope for them. Because if they want the lamp, we'd conflict, but if they don't, they can take other items. Alternatively, if I take the rope, which I value low, maybe that leaves the lamp for them if they value it.} \\
\\ 
\texttt{But my highest value is the lamp. So maybe I should go} \\
            }
        }
    }
    & \\ \\

    \theutterance \stepcounter{utterance}  
    & & & \multicolumn{4}{p{0.6\linewidth}}{
        \cellcolor[rgb]{0.9,0.9,0.9}{
            \makecell[{{p{\linewidth}}}]{
                \texttt{\tiny{[GM$\langle$P2]}}
                \texttt{<think>} \\
\texttt{Okay, let's try to figure out the best way to approach this. The available items are 1 lamp, 3 thimbles, and 1 rope. My values are lamp: 3, thimble: 2, rope: 1. The other player's values are secret, so I need to make a proposal that maximizes my points while allowing them to have a complementary set.} \\
\\ 
\texttt{First, I should consider what the other player might value. Since there's only one lamp and one rope, they might be high{-}value items for them too. But with three thimbles, maybe they can take some without conflicting. Let's think about possible scenarios.} \\
\\ 
\texttt{If I take the lamp (3 points), then the remaining items are 3 thimbles and 1 rope. The other player could take all 3 thimbles and the rope. If their value for thimbles is low, maybe they'd prefer the rope. But if they value thimbles more, taking all three might be good. However, since I don't know their values, I need to make a proposal that's flexible.} \\
\\ 
\texttt{Alternatively, if I take two thimbles (2*2=4) and the rope (1), that's 5 points. Then the other player could take the lamp and one thimble. But if they value the lamp highly, that might work. But if they also want thimbles, maybe they'd prefer more of them. However, there are only three, so if I take two, they can take one. But maybe they want the lamp and the rope instead.} \\
\\ 
\texttt{Another option is to take the lamp and one thimble (3+2=5). Then the other player could take two thimbles and the rope. But again, depends on their values.} \\
\\ 
\texttt{Wait, the key is to make a proposal that doesn't conflict with possible other proposals. Since the other player's values are unknown, I need to choose a proposal that allows them to have a set that doesn't overlap with mine. The safest way might be to take the lamp and leave the thimbles and rope for them. Because if they want the lamp, we'd conflict, but if they don't, they can take other items. Alternatively, if I take the rope, which I value low, maybe that leaves the lamp for them if they value it.} \\
\\ 
\texttt{But my highest value is the lamp. So maybe I should go} \\
            }
        }
    }
     \\ \\

    \theutterance \stepcounter{utterance}  
    & & & \multicolumn{2}{p{0.3\linewidth}}{
        \cellcolor[rgb]{0.9,0.9,0.9}{
            \makecell[{{p{\linewidth}}}]{
                \texttt{\tiny{[GM$|$GM]}}
                \texttt{continue} \\
            }
        }
    }
    & & \\ \\

    \theutterance \stepcounter{utterance}  
    & & & \multicolumn{2}{p{0.3\linewidth}}{
        \cellcolor[rgb]{0.9,0.9,0.9}{
            \makecell[{{p{\linewidth}}}]{
                \texttt{\tiny{[GM$|$GM]}}
                \texttt{<think>} \\
\texttt{Okay, let's try to figure out the best way to approach this. The available items are 1 lamp, 3 thimbles, and 1 rope. My values are lamp: 3, thimble: 2, rope: 1. The other player's values are secret, so I need to make a proposal that maximizes my points while allowing them to have a complementary set.} \\
\\ 
\texttt{First, I should consider what the other player might value. Since there's only one lamp and one rope, they might be high{-}value items for them too. But with three thimbles, maybe they can take some without conflicting. Let's think about possible scenarios.} \\
\\ 
\texttt{If I take the lamp (3 points), then the remaining items are 3 thimbles and 1 rope. The other player could take all 3 thimbles and the rope. If their value for thimbles is low, maybe they'd prefer the rope. But if they value thimbles more, taking all three might be good. However, since I don't know their values, I need to make a proposal that's flexible.} \\
\\ 
\texttt{Alternatively, if I take two thimbles (2*2=4) and the rope (1), that's 5 points. Then the other player could take the lamp and one thimble. But if they value the lamp highly, that might work. But if they also want thimbles, maybe they'd prefer more of them. However, there are only three, so if I take two, they can take one. But maybe they want the lamp and the rope instead.} \\
\\ 
\texttt{Another option is to take the lamp and one thimble (3+2=5). Then the other player could take two thimbles and the rope. But again, depends on their values.} \\
\\ 
\texttt{Wait, the key is to make a proposal that doesn't conflict with possible other proposals. Since the other player's values are unknown, I need to choose a proposal that allows them to have a set that doesn't overlap with mine. The safest way might be to take the lamp and leave the thimbles and rope for them. Because if they want the lamp, we'd conflict, but if they don't, they can take other items. Alternatively, if I take the rope, which I value low, maybe that leaves the lamp for them if they value it.} \\
\\ 
\texttt{But my highest value is the lamp. So maybe I should go} \\
            }
        }
    }
    & & \\ \\

    \theutterance \stepcounter{utterance}  
    & & \multicolumn{4}{p{0.6\linewidth}}{
        \cellcolor[rgb]{0.9,0.9,0.9}{
            \makecell[{{p{\linewidth}}}]{
                \texttt{\tiny{[P1$\langle$GM]}}
                \texttt{<think>} \\
\texttt{Okay, let's try to figure out the best way to approach this. The available items are 1 lamp, 3 thimbles, and 1 rope. My values are lamp: 3, thimble: 2, rope: 1. The other player's values are secret, so I need to make a proposal that maximizes my points while allowing them to have a complementary set.} \\
\\ 
\texttt{First, I should consider what the other player might value. Since there's only one lamp and one rope, they might be high{-}value items for them too. But with three thimbles, maybe they can take some without conflicting. Let's think about possible scenarios.} \\
\\ 
\texttt{If I take the lamp (3 points), then the remaining items are 3 thimbles and 1 rope. The other player could take all 3 thimbles and the rope. If their value for thimbles is low, maybe they'd prefer the rope. But if they value thimbles more, taking all three might be good. However, since I don't know their values, I need to make a proposal that's flexible.} \\
\\ 
\texttt{Alternatively, if I take two thimbles (2*2=4) and the rope (1), that's 5 points. Then the other player could take the lamp and one thimble. But if they value the lamp highly, that might work. But if they also want thimbles, maybe they'd prefer more of them. However, there are only three, so if I take two, they can take one. But maybe they want the lamp and the rope instead.} \\
\\ 
\texttt{Another option is to take the lamp and one thimble (3+2=5). Then the other player could take two thimbles and the rope. But again, depends on their values.} \\
\\ 
\texttt{Wait, the key is to make a proposal that doesn't conflict with possible other proposals. Since the other player's values are unknown, I need to choose a proposal that allows them to have a set that doesn't overlap with mine. The safest way might be to take the lamp and leave the thimbles and rope for them. Because if they want the lamp, we'd conflict, but if they don't, they can take other items. Alternatively, if I take the rope, which I value low, maybe that leaves the lamp for them if they value it.} \\
\\ 
\texttt{But my highest value is the lamp. So maybe I should go} \\
            }
        }
    }
    & \\ \\

    \theutterance \stepcounter{utterance}  
    & \multicolumn{4}{p{0.6\linewidth}}{
        \cellcolor[rgb]{0.9,0.9,0.9}{
            \makecell[{{p{\linewidth}}}]{
                \texttt{\tiny{[P1$\rangle$GM]}}
                \texttt{<think>} \\
\texttt{Okay, let's try to figure out the best way to approach this. The available items are 1 lamp, 3 thimbles, and 1 rope. My values are lamp: 3, thimble: 2, rope: 1. The other player's values are secret, so I need to make a proposal that maximizes my points while allowing them to have a complementary set.} \\
\\ 
\texttt{First, I should consider what the other player might value. Since there's only one lamp and one rope, they might be high{-}value items for them too. But with three thimbles, maybe they can take some without conflicting. Let's think about possible scenarios.} \\
\\ 
\texttt{If I take the lamp (3 points), then the remaining items are 3 thimbles and 1 rope. The other player could take all 3 thimbles and the rope. If their value for thimbles is low, maybe they'd prefer the rope. But if they value thimbles more, taking all three might be good. However, since I don't know their values, I need to make a proposal that's flexible.} \\
\\ 
\texttt{Alternatively, if I take two thimbles (2*2=4) and the rope (1), that's 5 points. Then the other player could take the lamp and one thimble. But if they value the lamp highly, that might work. But if they also want thimbles, maybe they'd prefer more of them. However, there are only three, so if I take two, they can take one. But maybe they want the lamp and the rope instead.} \\
\\ 
\texttt{Another option is to take the lamp and one thimble (3+2=5). Then the other player could take two thimbles and the rope. But again, depends on their values.} \\
\\ 
\texttt{Wait, the key is to make a proposal that doesn't conflict with possible other proposals. Since the other player's values are unknown, I need to choose a proposal that allows them to have a set that doesn't overlap with mine. The safest way might be to take the lamp and leave the thimbles and rope for them. Because if they want the lamp, we'd conflict, but if they don't, they can take other items. Alternatively, if I take the rope, which I value low, maybe that leaves the lamp for them if they value it.} \\
\\ 
\texttt{But my highest value is the lamp. So maybe I should go} \\
            }
        }
    }
    & & \\ \\

    \theutterance \stepcounter{utterance}  
    & & & \multicolumn{2}{p{0.3\linewidth}}{
        \cellcolor[rgb]{0.9,0.9,0.9}{
            \makecell[{{p{\linewidth}}}]{
                \texttt{\tiny{[GM$|$GM]}}
                \texttt{continue} \\
            }
        }
    }
    & & \\ \\

    \theutterance \stepcounter{utterance}  
    & & & \multicolumn{2}{p{0.3\linewidth}}{
        \cellcolor[rgb]{0.9,0.9,0.9}{
            \makecell[{{p{\linewidth}}}]{
                \texttt{\tiny{[GM$|$GM]}}
                \texttt{<think>} \\
\texttt{Okay, let's try to figure out the best way to approach this. The available items are 1 lamp, 3 thimbles, and 1 rope. My values are lamp: 3, thimble: 2, rope: 1. The other player's values are secret, so I need to make a proposal that maximizes my points while allowing them to have a complementary set.} \\
\\ 
\texttt{First, I should consider what the other player might value. Since there's only one lamp and one rope, they might be high{-}value items for them too. But with three thimbles, maybe they can take some without conflicting. Let's think about possible scenarios.} \\
\\ 
\texttt{If I take the lamp (3 points), then the remaining items are 3 thimbles and 1 rope. The other player could take all 3 thimbles and the rope. If their value for thimbles is low, maybe they'd prefer the rope. But if they value thimbles more, taking all three might be good. However, since I don't know their values, I need to make a proposal that's flexible.} \\
\\ 
\texttt{Alternatively, if I take two thimbles (2*2=4) and the rope (1), that's 5 points. Then the other player could take the lamp and one thimble. But if they value the lamp highly, that might work. But if they also want thimbles, maybe they'd prefer more of them. However, there are only three, so if I take two, they can take one. But maybe they want the lamp and the rope instead.} \\
\\ 
\texttt{Another option is to take the lamp and one thimble (3+2=5). Then the other player could take two thimbles and the rope. But again, depends on their values.} \\
\\ 
\texttt{Wait, the key is to make a proposal that doesn't conflict with possible other proposals. Since the other player's values are unknown, I need to choose a proposal that allows them to have a set that doesn't overlap with mine. The safest way might be to take the lamp and leave the thimbles and rope for them. Because if they want the lamp, we'd conflict, but if they don't, they can take other items. Alternatively, if I take the rope, which I value low, maybe that leaves the lamp for them if they value it.} \\
\\ 
\texttt{But my highest value is the lamp. So maybe I should go} \\
            }
        }
    }
    & & \\ \\

    \theutterance \stepcounter{utterance}  
    & & \multicolumn{4}{p{0.6\linewidth}}{
        \cellcolor[rgb]{0.9,0.9,0.9}{
            \makecell[{{p{\linewidth}}}]{
                \texttt{\tiny{[GM$\rangle$P2]}}
                \texttt{<think>} \\
\texttt{Okay, let's try to figure out the best way to approach this. The available items are 1 lamp, 3 thimbles, and 1 rope. My values are lamp: 3, thimble: 2, rope: 1. The other player's values are secret, so I need to make a proposal that maximizes my points while allowing them to have a complementary set.} \\
\\ 
\texttt{First, I should consider what the other player might value. Since there's only one lamp and one rope, they might be high{-}value items for them too. But with three thimbles, maybe they can take some without conflicting. Let's think about possible scenarios.} \\
\\ 
\texttt{If I take the lamp (3 points), then the remaining items are 3 thimbles and 1 rope. The other player could take all 3 thimbles and the rope. If their value for thimbles is low, maybe they'd prefer the rope. But if they value thimbles more, taking all three might be good. However, since I don't know their values, I need to make a proposal that's flexible.} \\
\\ 
\texttt{Alternatively, if I take two thimbles (2*2=4) and the rope (1), that's 5 points. Then the other player could take the lamp and one thimble. But if they value the lamp highly, that might work. But if they also want thimbles, maybe they'd prefer more of them. However, there are only three, so if I take two, they can take one. But maybe they want the lamp and the rope instead.} \\
\\ 
\texttt{Another option is to take the lamp and one thimble (3+2=5). Then the other player could take two thimbles and the rope. But again, depends on their values.} \\
\\ 
\texttt{Wait, the key is to make a proposal that doesn't conflict with possible other proposals. Since the other player's values are unknown, I need to choose a proposal that allows them to have a set that doesn't overlap with mine. The safest way might be to take the lamp and leave the thimbles and rope for them. Because if they want the lamp, we'd conflict, but if they don't, they can take other items. Alternatively, if I take the rope, which I value low, maybe that leaves the lamp for them if they value it.} \\
\\ 
\texttt{But my highest value is the lamp. So maybe I should go} \\
            }
        }
    }
    & \\ \\

    \theutterance \stepcounter{utterance}  
    & & & \multicolumn{4}{p{0.6\linewidth}}{
        \cellcolor[rgb]{0.9,0.9,0.9}{
            \makecell[{{p{\linewidth}}}]{
                \texttt{\tiny{[GM$\langle$P2]}}
                \texttt{<think>} \\
\texttt{Okay, let's try to figure out the best way to approach this. The available items are 1 lamp, 3 thimbles, and 1 rope. My values are lamp: 3, thimble: 2, rope: 1. The other player's values are secret, so I need to make a proposal that maximizes my points while allowing them to have a complementary set.} \\
\\ 
\texttt{First, I should consider what the other player might value. Since there's only one lamp and one rope, they might be high{-}value items for them too. But with three thimbles, maybe they can take some without conflicting. Let's think about possible scenarios.} \\
\\ 
\texttt{If I take the lamp (3 points), then the remaining items are 3 thimbles and 1 rope. The other player could take all 3 thimbles and the rope. If their value for thimbles is low, maybe they'd prefer the rope. But if they value thimbles more, taking all three might be good. However, since I don't know their values, I need to make a proposal that's flexible.} \\
\\ 
\texttt{Alternatively, if I take two thimbles (2*2=4) and the rope (1), that's 5 points. Then the other player could take the lamp and one thimble. But if they value the lamp highly, that might work. But if they also want thimbles, maybe they'd prefer more of them. However, there are only three, so if I take two, they can take one. But maybe they want the lamp and the rope instead.} \\
\\ 
\texttt{Another option is to take the lamp and one thimble (3+2=5). Then the other player could take two thimbles and the rope. But again, depends on their values.} \\
\\ 
\texttt{Wait, the key is to make a proposal that doesn't conflict with possible other proposals. Since the other player's values are unknown, I need to choose a proposal that allows them to have a set that doesn't overlap with mine. The safest way might be to take the lamp and leave the thimbles and rope for them. Because if they want the lamp, we'd conflict, but if they don't, they can take other items. Alternatively, if I take the rope, which I value low, maybe that leaves the lamp for them if they value it.} \\
\\ 
\texttt{But my highest value is the lamp. So maybe I should go} \\
            }
        }
    }
     \\ \\

    \theutterance \stepcounter{utterance}  
    & & & \multicolumn{2}{p{0.3\linewidth}}{
        \cellcolor[rgb]{0.9,0.9,0.9}{
            \makecell[{{p{\linewidth}}}]{
                \texttt{\tiny{[GM$|$GM]}}
                \texttt{continue} \\
            }
        }
    }
    & & \\ \\

    \theutterance \stepcounter{utterance}  
    & & & \multicolumn{2}{p{0.3\linewidth}}{
        \cellcolor[rgb]{0.9,0.9,0.9}{
            \makecell[{{p{\linewidth}}}]{
                \texttt{\tiny{[GM$|$GM]}}
                \texttt{<think>} \\
\texttt{Okay, let's try to figure out the best way to approach this. The available items are 1 lamp, 3 thimbles, and 1 rope. My values are lamp: 3, thimble: 2, rope: 1. The other player's values are secret, so I need to make a proposal that maximizes my points while allowing them to have a complementary set.} \\
\\ 
\texttt{First, I should consider what the other player might value. Since there's only one lamp and one rope, they might be high{-}value items for them too. But with three thimbles, maybe they can take some without conflicting. Let's think about possible scenarios.} \\
\\ 
\texttt{If I take the lamp (3 points), then the remaining items are 3 thimbles and 1 rope. The other player could take all 3 thimbles and the rope. If their value for thimbles is low, maybe they'd prefer the rope. But if they value thimbles more, taking all three might be good. However, since I don't know their values, I need to make a proposal that's flexible.} \\
\\ 
\texttt{Alternatively, if I take two thimbles (2*2=4) and the rope (1), that's 5 points. Then the other player could take the lamp and one thimble. But if they value the lamp highly, that might work. But if they also want thimbles, maybe they'd prefer more of them. However, there are only three, so if I take two, they can take one. But maybe they want the lamp and the rope instead.} \\
\\ 
\texttt{Another option is to take the lamp and one thimble (3+2=5). Then the other player could take two thimbles and the rope. But again, depends on their values.} \\
\\ 
\texttt{Wait, the key is to make a proposal that doesn't conflict with possible other proposals. Since the other player's values are unknown, I need to choose a proposal that allows them to have a set that doesn't overlap with mine. The safest way might be to take the lamp and leave the thimbles and rope for them. Because if they want the lamp, we'd conflict, but if they don't, they can take other items. Alternatively, if I take the rope, which I value low, maybe that leaves the lamp for them if they value it.} \\
\\ 
\texttt{But my highest value is the lamp. So maybe I should go} \\
            }
        }
    }
    & & \\ \\

    \theutterance \stepcounter{utterance}  
    & & \multicolumn{4}{p{0.6\linewidth}}{
        \cellcolor[rgb]{0.9,0.9,0.9}{
            \makecell[{{p{\linewidth}}}]{
                \texttt{\tiny{[P1$\langle$GM]}}
                \texttt{<think>} \\
\texttt{Okay, let's try to figure out the best way to approach this. The available items are 1 lamp, 3 thimbles, and 1 rope. My values are lamp: 3, thimble: 2, rope: 1. The other player's values are secret, so I need to make a proposal that maximizes my points while allowing them to have a complementary set.} \\
\\ 
\texttt{First, I should consider what the other player might value. Since there's only one lamp and one rope, they might be high{-}value items for them too. But with three thimbles, maybe they can take some without conflicting. Let's think about possible scenarios.} \\
\\ 
\texttt{If I take the lamp (3 points), then the remaining items are 3 thimbles and 1 rope. The other player could take all 3 thimbles and the rope. If their value for thimbles is low, maybe they'd prefer the rope. But if they value thimbles more, taking all three might be good. However, since I don't know their values, I need to make a proposal that's flexible.} \\
\\ 
\texttt{Alternatively, if I take two thimbles (2*2=4) and the rope (1), that's 5 points. Then the other player could take the lamp and one thimble. But if they value the lamp highly, that might work. But if they also want thimbles, maybe they'd prefer more of them. However, there are only three, so if I take two, they can take one. But maybe they want the lamp and the rope instead.} \\
\\ 
\texttt{Another option is to take the lamp and one thimble (3+2=5). Then the other player could take two thimbles and the rope. But again, depends on their values.} \\
\\ 
\texttt{Wait, the key is to make a proposal that doesn't conflict with possible other proposals. Since the other player's values are unknown, I need to choose a proposal that allows them to have a set that doesn't overlap with mine. The safest way might be to take the lamp and leave the thimbles and rope for them. Because if they want the lamp, we'd conflict, but if they don't, they can take other items. Alternatively, if I take the rope, which I value low, maybe that leaves the lamp for them if they value it.} \\
\\ 
\texttt{But my highest value is the lamp. So maybe I should go} \\
            }
        }
    }
    & \\ \\

    \theutterance \stepcounter{utterance}  
    & \multicolumn{4}{p{0.6\linewidth}}{
        \cellcolor[rgb]{0.9,0.9,0.9}{
            \makecell[{{p{\linewidth}}}]{
                \texttt{\tiny{[P1$\rangle$GM]}}
                \texttt{<think>} \\
\texttt{Okay, let's try to figure out the best way to approach this. The available items are 1 lamp, 3 thimbles, and 1 rope. My values are lamp: 3, thimble: 2, rope: 1. The other player's values are secret, so I need to make a proposal that maximizes my points while allowing them to have a complementary set.} \\
\\ 
\texttt{First, I should consider what the other player might value. Since there's only one lamp and one rope, they might be high{-}value items for them too. But with three thimbles, maybe they can take some without conflicting. Let's think about possible scenarios.} \\
\\ 
\texttt{If I take the lamp (3 points), then the remaining items are 3 thimbles and 1 rope. The other player could take all 3 thimbles and the rope. If their value for thimbles is low, maybe they'd prefer the rope. But if they value thimbles more, taking all three might be good. However, since I don't know their values, I need to make a proposal that's flexible.} \\
\\ 
\texttt{Alternatively, if I take two thimbles (2*2=4) and the rope (1), that's 5 points. Then the other player could take the lamp and one thimble. But if they value the lamp highly, that might work. But if they also want thimbles, maybe they'd prefer more of them. However, there are only three, so if I take two, they can take one. But maybe they want the lamp and the rope instead.} \\
\\ 
\texttt{Another option is to take the lamp and one thimble (3+2=5). Then the other player could take two thimbles and the rope. But again, depends on their values.} \\
\\ 
\texttt{Wait, the key is to make a proposal that doesn't conflict with possible other proposals. Since the other player's values are unknown, I need to choose a proposal that allows them to have a set that doesn't overlap with mine. The safest way might be to take the lamp and leave the thimbles and rope for them. Because if they want the lamp, we'd conflict, but if they don't, they can take other items. Alternatively, if I take the rope, which I value low, maybe that leaves the lamp for them if they value it.} \\
\\ 
\texttt{But my highest value is the lamp. So maybe I should go} \\
            }
        }
    }
    & & \\ \\

    \theutterance \stepcounter{utterance}  
    & & & \multicolumn{2}{p{0.3\linewidth}}{
        \cellcolor[rgb]{0.9,0.9,0.9}{
            \makecell[{{p{\linewidth}}}]{
                \texttt{\tiny{[GM$|$GM]}}
                \texttt{continue} \\
            }
        }
    }
    & & \\ \\

    \theutterance \stepcounter{utterance}  
    & & & \multicolumn{2}{p{0.3\linewidth}}{
        \cellcolor[rgb]{0.9,0.9,0.9}{
            \makecell[{{p{\linewidth}}}]{
                \texttt{\tiny{[GM$|$GM]}}
                \texttt{<think>} \\
\texttt{Okay, let's try to figure out the best way to approach this. The available items are 1 lamp, 3 thimbles, and 1 rope. My values are lamp: 3, thimble: 2, rope: 1. The other player's values are secret, so I need to make a proposal that maximizes my points while allowing them to have a complementary set.} \\
\\ 
\texttt{First, I should consider what the other player might value. Since there's only one lamp and one rope, they might be high{-}value items for them too. But with three thimbles, maybe they can take some without conflicting. Let's think about possible scenarios.} \\
\\ 
\texttt{If I take the lamp (3 points), then the remaining items are 3 thimbles and 1 rope. The other player could take all 3 thimbles and the rope. If their value for thimbles is low, maybe they'd prefer the rope. But if they value thimbles more, taking all three might be good. However, since I don't know their values, I need to make a proposal that's flexible.} \\
\\ 
\texttt{Alternatively, if I take two thimbles (2*2=4) and the rope (1), that's 5 points. Then the other player could take the lamp and one thimble. But if they value the lamp highly, that might work. But if they also want thimbles, maybe they'd prefer more of them. However, there are only three, so if I take two, they can take one. But maybe they want the lamp and the rope instead.} \\
\\ 
\texttt{Another option is to take the lamp and one thimble (3+2=5). Then the other player could take two thimbles and the rope. But again, depends on their values.} \\
\\ 
\texttt{Wait, the key is to make a proposal that doesn't conflict with possible other proposals. Since the other player's values are unknown, I need to choose a proposal that allows them to have a set that doesn't overlap with mine. The safest way might be to take the lamp and leave the thimbles and rope for them. Because if they want the lamp, we'd conflict, but if they don't, they can take other items. Alternatively, if I take the rope, which I value low, maybe that leaves the lamp for them if they value it.} \\
\\ 
\texttt{But my highest value is the lamp. So maybe I should go} \\
            }
        }
    }
    & & \\ \\

    \theutterance \stepcounter{utterance}  
    & & \multicolumn{4}{p{0.6\linewidth}}{
        \cellcolor[rgb]{0.9,0.9,0.9}{
            \makecell[{{p{\linewidth}}}]{
                \texttt{\tiny{[GM$\rangle$P2]}}
                \texttt{<think>} \\
\texttt{Okay, let's try to figure out the best way to approach this. The available items are 1 lamp, 3 thimbles, and 1 rope. My values are lamp: 3, thimble: 2, rope: 1. The other player's values are secret, so I need to make a proposal that maximizes my points while allowing them to have a complementary set.} \\
\\ 
\texttt{First, I should consider what the other player might value. Since there's only one lamp and one rope, they might be high{-}value items for them too. But with three thimbles, maybe they can take some without conflicting. Let's think about possible scenarios.} \\
\\ 
\texttt{If I take the lamp (3 points), then the remaining items are 3 thimbles and 1 rope. The other player could take all 3 thimbles and the rope. If their value for thimbles is low, maybe they'd prefer the rope. But if they value thimbles more, taking all three might be good. However, since I don't know their values, I need to make a proposal that's flexible.} \\
\\ 
\texttt{Alternatively, if I take two thimbles (2*2=4) and the rope (1), that's 5 points. Then the other player could take the lamp and one thimble. But if they value the lamp highly, that might work. But if they also want thimbles, maybe they'd prefer more of them. However, there are only three, so if I take two, they can take one. But maybe they want the lamp and the rope instead.} \\
\\ 
\texttt{Another option is to take the lamp and one thimble (3+2=5). Then the other player could take two thimbles and the rope. But again, depends on their values.} \\
\\ 
\texttt{Wait, the key is to make a proposal that doesn't conflict with possible other proposals. Since the other player's values are unknown, I need to choose a proposal that allows them to have a set that doesn't overlap with mine. The safest way might be to take the lamp and leave the thimbles and rope for them. Because if they want the lamp, we'd conflict, but if they don't, they can take other items. Alternatively, if I take the rope, which I value low, maybe that leaves the lamp for them if they value it.} \\
\\ 
\texttt{But my highest value is the lamp. So maybe I should go} \\
            }
        }
    }
    & \\ \\

    \theutterance \stepcounter{utterance}  
    & & & \multicolumn{4}{p{0.6\linewidth}}{
        \cellcolor[rgb]{0.9,0.9,0.9}{
            \makecell[{{p{\linewidth}}}]{
                \texttt{\tiny{[GM$\langle$P2]}}
                \texttt{<think>} \\
\texttt{Okay, let's try to figure out the best way to approach this. The available items are 1 lamp, 3 thimbles, and 1 rope. My values are lamp: 3, thimble: 2, rope: 1. The other player's values are secret, so I need to make a proposal that maximizes my points while allowing them to have a complementary set.} \\
\\ 
\texttt{First, I should consider what the other player might value. Since there's only one lamp and one rope, they might be high{-}value items for them too. But with three thimbles, maybe they can take some without conflicting. Let's think about possible scenarios.} \\
\\ 
\texttt{If I take the lamp (3 points), then the remaining items are 3 thimbles and 1 rope. The other player could take all 3 thimbles and the rope. If their value for thimbles is low, maybe they'd prefer the rope. But if they value thimbles more, taking all three might be good. However, since I don't know their values, I need to make a proposal that's flexible.} \\
\\ 
\texttt{Alternatively, if I take two thimbles (2*2=4) and the rope (1), that's 5 points. Then the other player could take the lamp and one thimble. But if they value the lamp highly, that might work. But if they also want thimbles, maybe they'd prefer more of them. However, there are only three, so if I take two, they can take one. But maybe they want the lamp and the rope instead.} \\
\\ 
\texttt{Another option is to take the lamp and one thimble (3+2=5). Then the other player could take two thimbles and the rope. But again, depends on their values.} \\
\\ 
\texttt{Wait, the key is to make a proposal that doesn't conflict with possible other proposals. Since the other player's values are unknown, I need to choose a proposal that allows them to have a set that doesn't overlap with mine. The safest way might be to take the lamp and leave the thimbles and rope for them. Because if they want the lamp, we'd conflict, but if they don't, they can take other items. Alternatively, if I take the rope, which I value low, maybe that leaves the lamp for them if they value it.} \\
\\ 
\texttt{But my highest value is the lamp. So maybe I should go} \\
            }
        }
    }
     \\ \\

    \theutterance \stepcounter{utterance}  
    & & & \multicolumn{2}{p{0.3\linewidth}}{
        \cellcolor[rgb]{0.9,0.9,0.9}{
            \makecell[{{p{\linewidth}}}]{
                \texttt{\tiny{[GM$|$GM]}}
                \texttt{continue} \\
            }
        }
    }
    & & \\ \\

    \theutterance \stepcounter{utterance}  
    & & & \multicolumn{2}{p{0.3\linewidth}}{
        \cellcolor[rgb]{0.9,0.9,0.9}{
            \makecell[{{p{\linewidth}}}]{
                \texttt{\tiny{[GM$|$GM]}}
                \texttt{<think>} \\
\texttt{Okay, let's try to figure out the best way to approach this. The available items are 1 lamp, 3 thimbles, and 1 rope. My values are lamp: 3, thimble: 2, rope: 1. The other player's values are secret, so I need to make a proposal that maximizes my points while allowing them to have a complementary set.} \\
\\ 
\texttt{First, I should consider what the other player might value. Since there's only one lamp and one rope, they might be high{-}value items for them too. But with three thimbles, maybe they can take some without conflicting. Let's think about possible scenarios.} \\
\\ 
\texttt{If I take the lamp (3 points), then the remaining items are 3 thimbles and 1 rope. The other player could take all 3 thimbles and the rope. If their value for thimbles is low, maybe they'd prefer the rope. But if they value thimbles more, taking all three might be good. However, since I don't know their values, I need to make a proposal that's flexible.} \\
\\ 
\texttt{Alternatively, if I take two thimbles (2*2=4) and the rope (1), that's 5 points. Then the other player could take the lamp and one thimble. But if they value the lamp highly, that might work. But if they also want thimbles, maybe they'd prefer more of them. However, there are only three, so if I take two, they can take one. But maybe they want the lamp and the rope instead.} \\
\\ 
\texttt{Another option is to take the lamp and one thimble (3+2=5). Then the other player could take two thimbles and the rope. But again, depends on their values.} \\
\\ 
\texttt{Wait, the key is to make a proposal that doesn't conflict with possible other proposals. Since the other player's values are unknown, I need to choose a proposal that allows them to have a set that doesn't overlap with mine. The safest way might be to take the lamp and leave the thimbles and rope for them. Because if they want the lamp, we'd conflict, but if they don't, they can take other items. Alternatively, if I take the rope, which I value low, maybe that leaves the lamp for them if they value it.} \\
\\ 
\texttt{But my highest value is the lamp. So maybe I should go} \\
            }
        }
    }
    & & \\ \\

    \theutterance \stepcounter{utterance}  
    & & \multicolumn{4}{p{0.6\linewidth}}{
        \cellcolor[rgb]{0.9,0.9,0.9}{
            \makecell[{{p{\linewidth}}}]{
                \texttt{\tiny{[P1$\langle$GM]}}
                \texttt{<think>} \\
\texttt{Okay, let's try to figure out the best way to approach this. The available items are 1 lamp, 3 thimbles, and 1 rope. My values are lamp: 3, thimble: 2, rope: 1. The other player's values are secret, so I need to make a proposal that maximizes my points while allowing them to have a complementary set.} \\
\\ 
\texttt{First, I should consider what the other player might value. Since there's only one lamp and one rope, they might be high{-}value items for them too. But with three thimbles, maybe they can take some without conflicting. Let's think about possible scenarios.} \\
\\ 
\texttt{If I take the lamp (3 points), then the remaining items are 3 thimbles and 1 rope. The other player could take all 3 thimbles and the rope. If their value for thimbles is low, maybe they'd prefer the rope. But if they value thimbles more, taking all three might be good. However, since I don't know their values, I need to make a proposal that's flexible.} \\
\\ 
\texttt{Alternatively, if I take two thimbles (2*2=4) and the rope (1), that's 5 points. Then the other player could take the lamp and one thimble. But if they value the lamp highly, that might work. But if they also want thimbles, maybe they'd prefer more of them. However, there are only three, so if I take two, they can take one. But maybe they want the lamp and the rope instead.} \\
\\ 
\texttt{Another option is to take the lamp and one thimble (3+2=5). Then the other player could take two thimbles and the rope. But again, depends on their values.} \\
\\ 
\texttt{Wait, the key is to make a proposal that doesn't conflict with possible other proposals. Since the other player's values are unknown, I need to choose a proposal that allows them to have a set that doesn't overlap with mine. The safest way might be to take the lamp and leave the thimbles and rope for them. Because if they want the lamp, we'd conflict, but if they don't, they can take other items. Alternatively, if I take the rope, which I value low, maybe that leaves the lamp for them if they value it.} \\
\\ 
\texttt{But my highest value is the lamp. So maybe I should go} \\
            }
        }
    }
    & \\ \\

    \theutterance \stepcounter{utterance}  
    & \multicolumn{4}{p{0.6\linewidth}}{
        \cellcolor[rgb]{0.9,0.9,0.9}{
            \makecell[{{p{\linewidth}}}]{
                \texttt{\tiny{[P1$\rangle$GM]}}
                \texttt{<think>} \\
\texttt{Okay, let's try to figure out the best way to approach this. The available items are 1 lamp, 3 thimbles, and 1 rope. My values are lamp: 3, thimble: 2, rope: 1. The other player's values are secret, so I need to make a proposal that maximizes my points while allowing them to have a complementary set.} \\
\\ 
\texttt{First, I should consider what the other player might value. Since there's only one lamp and one rope, they might be high{-}value items for them too. But with three thimbles, maybe they can take some without conflicting. Let's think about possible scenarios.} \\
\\ 
\texttt{If I take the lamp (3 points), then the remaining items are 3 thimbles and 1 rope. The other player could take all 3 thimbles and the rope. If their value for thimbles is low, maybe they'd prefer the rope. But if they value thimbles more, taking all three might be good. However, since I don't know their values, I need to make a proposal that's flexible.} \\
\\ 
\texttt{Alternatively, if I take two thimbles (2*2=4) and the rope (1), that's 5 points. Then the other player could take the lamp and one thimble. But if they value the lamp highly, that might work. But if they also want thimbles, maybe they'd prefer more of them. However, there are only three, so if I take two, they can take one. But maybe they want the lamp and the rope instead.} \\
\\ 
\texttt{Another option is to take the lamp and one thimble (3+2=5). Then the other player could take two thimbles and the rope. But again, depends on their values.} \\
\\ 
\texttt{Wait, the key is to make a proposal that doesn't conflict with possible other proposals. Since the other player's values are unknown, I need to choose a proposal that allows them to have a set that doesn't overlap with mine. The safest way might be to take the lamp and leave the thimbles and rope for them. Because if they want the lamp, we'd conflict, but if they don't, they can take other items. Alternatively, if I take the rope, which I value low, maybe that leaves the lamp for them if they value it.} \\
\\ 
\texttt{But my highest value is the lamp. So maybe I should go} \\
            }
        }
    }
    & & \\ \\

    \theutterance \stepcounter{utterance}  
    & & & \multicolumn{2}{p{0.3\linewidth}}{
        \cellcolor[rgb]{0.9,0.9,0.9}{
            \makecell[{{p{\linewidth}}}]{
                \texttt{\tiny{[GM$|$GM]}}
                \texttt{continue} \\
            }
        }
    }
    & & \\ \\

    \theutterance \stepcounter{utterance}  
    & & & \multicolumn{2}{p{0.3\linewidth}}{
        \cellcolor[rgb]{0.9,0.9,0.9}{
            \makecell[{{p{\linewidth}}}]{
                \texttt{\tiny{[GM$|$GM]}}
                \texttt{<think>} \\
\texttt{Okay, let's try to figure out the best way to approach this. The available items are 1 lamp, 3 thimbles, and 1 rope. My values are lamp: 3, thimble: 2, rope: 1. The other player's values are secret, so I need to make a proposal that maximizes my points while allowing them to have a complementary set.} \\
\\ 
\texttt{First, I should consider what the other player might value. Since there's only one lamp and one rope, they might be high{-}value items for them too. But with three thimbles, maybe they can take some without conflicting. Let's think about possible scenarios.} \\
\\ 
\texttt{If I take the lamp (3 points), then the remaining items are 3 thimbles and 1 rope. The other player could take all 3 thimbles and the rope. If their value for thimbles is low, maybe they'd prefer the rope. But if they value thimbles more, taking all three might be good. However, since I don't know their values, I need to make a proposal that's flexible.} \\
\\ 
\texttt{Alternatively, if I take two thimbles (2*2=4) and the rope (1), that's 5 points. Then the other player could take the lamp and one thimble. But if they value the lamp highly, that might work. But if they also want thimbles, maybe they'd prefer more of them. However, there are only three, so if I take two, they can take one. But maybe they want the lamp and the rope instead.} \\
\\ 
\texttt{Another option is to take the lamp and one thimble (3+2=5). Then the other player could take two thimbles and the rope. But again, depends on their values.} \\
\\ 
\texttt{Wait, the key is to make a proposal that doesn't conflict with possible other proposals. Since the other player's values are unknown, I need to choose a proposal that allows them to have a set that doesn't overlap with mine. The safest way might be to take the lamp and leave the thimbles and rope for them. Because if they want the lamp, we'd conflict, but if they don't, they can take other items. Alternatively, if I take the rope, which I value low, maybe that leaves the lamp for them if they value it.} \\
\\ 
\texttt{But my highest value is the lamp. So maybe I should go} \\
            }
        }
    }
    & & \\ \\

    \theutterance \stepcounter{utterance}  
    & & \multicolumn{4}{p{0.6\linewidth}}{
        \cellcolor[rgb]{0.9,0.9,0.9}{
            \makecell[{{p{\linewidth}}}]{
                \texttt{\tiny{[GM$\rangle$P2]}}
                \texttt{<think>} \\
\texttt{Okay, let's try to figure out the best way to approach this. The available items are 1 lamp, 3 thimbles, and 1 rope. My values are lamp: 3, thimble: 2, rope: 1. The other player's values are secret, so I need to make a proposal that maximizes my points while allowing them to have a complementary set.} \\
\\ 
\texttt{First, I should consider what the other player might value. Since there's only one lamp and one rope, they might be high{-}value items for them too. But with three thimbles, maybe they can take some without conflicting. Let's think about possible scenarios.} \\
\\ 
\texttt{If I take the lamp (3 points), then the remaining items are 3 thimbles and 1 rope. The other player could take all 3 thimbles and the rope. If their value for thimbles is low, maybe they'd prefer the rope. But if they value thimbles more, taking all three might be good. However, since I don't know their values, I need to make a proposal that's flexible.} \\
\\ 
\texttt{Alternatively, if I take two thimbles (2*2=4) and the rope (1), that's 5 points. Then the other player could take the lamp and one thimble. But if they value the lamp highly, that might work. But if they also want thimbles, maybe they'd prefer more of them. However, there are only three, so if I take two, they can take one. But maybe they want the lamp and the rope instead.} \\
\\ 
\texttt{Another option is to take the lamp and one thimble (3+2=5). Then the other player could take two thimbles and the rope. But again, depends on their values.} \\
\\ 
\texttt{Wait, the key is to make a proposal that doesn't conflict with possible other proposals. Since the other player's values are unknown, I need to choose a proposal that allows them to have a set that doesn't overlap with mine. The safest way might be to take the lamp and leave the thimbles and rope for them. Because if they want the lamp, we'd conflict, but if they don't, they can take other items. Alternatively, if I take the rope, which I value low, maybe that leaves the lamp for them if they value it.} \\
\\ 
\texttt{But my highest value is the lamp. So maybe I should go} \\
            }
        }
    }
    & \\ \\

    \theutterance \stepcounter{utterance}  
    & & & \multicolumn{4}{p{0.6\linewidth}}{
        \cellcolor[rgb]{0.9,0.9,0.9}{
            \makecell[{{p{\linewidth}}}]{
                \texttt{\tiny{[GM$\langle$P2]}}
                \texttt{<think>} \\
\texttt{Okay, let's try to figure out the best way to approach this. The available items are 1 lamp, 3 thimbles, and 1 rope. My values are lamp: 3, thimble: 2, rope: 1. The other player's values are secret, so I need to make a proposal that maximizes my points while allowing them to have a complementary set.} \\
\\ 
\texttt{First, I should consider what the other player might value. Since there's only one lamp and one rope, they might be high{-}value items for them too. But with three thimbles, maybe they can take some without conflicting. Let's think about possible scenarios.} \\
\\ 
\texttt{If I take the lamp (3 points), then the remaining items are 3 thimbles and 1 rope. The other player could take all 3 thimbles and the rope. If their value for thimbles is low, maybe they'd prefer the rope. But if they value thimbles more, taking all three might be good. However, since I don't know their values, I need to make a proposal that's flexible.} \\
\\ 
\texttt{Alternatively, if I take two thimbles (2*2=4) and the rope (1), that's 5 points. Then the other player could take the lamp and one thimble. But if they value the lamp highly, that might work. But if they also want thimbles, maybe they'd prefer more of them. However, there are only three, so if I take two, they can take one. But maybe they want the lamp and the rope instead.} \\
\\ 
\texttt{Another option is to take the lamp and one thimble (3+2=5). Then the other player could take two thimbles and the rope. But again, depends on their values.} \\
\\ 
\texttt{Wait, the key is to make a proposal that doesn't conflict with possible other proposals. Since the other player's values are unknown, I need to choose a proposal that allows them to have a set that doesn't overlap with mine. The safest way might be to take the lamp and leave the thimbles and rope for them. Because if they want the lamp, we'd conflict, but if they don't, they can take other items. Alternatively, if I take the rope, which I value low, maybe that leaves the lamp for them if they value it.} \\
\\ 
\texttt{But my highest value is the lamp. So maybe I should go} \\
            }
        }
    }
     \\ \\

    \theutterance \stepcounter{utterance}  
    & & & \multicolumn{2}{p{0.3\linewidth}}{
        \cellcolor[rgb]{0.9,0.9,0.9}{
            \makecell[{{p{\linewidth}}}]{
                \texttt{\tiny{[GM$|$GM]}}
                \texttt{continue} \\
            }
        }
    }
    & & \\ \\

    \theutterance \stepcounter{utterance}  
    & & & \multicolumn{2}{p{0.3\linewidth}}{
        \cellcolor[rgb]{0.9,0.9,0.9}{
            \makecell[{{p{\linewidth}}}]{
                \texttt{\tiny{[GM$|$GM]}}
                \texttt{<think>} \\
\texttt{Okay, let's try to figure out the best way to approach this. The available items are 1 lamp, 3 thimbles, and 1 rope. My values are lamp: 3, thimble: 2, rope: 1. The other player's values are secret, so I need to make a proposal that maximizes my points while allowing them to have a complementary set.} \\
\\ 
\texttt{First, I should consider what the other player might value. Since there's only one lamp and one rope, they might be high{-}value items for them too. But with three thimbles, maybe they can take some without conflicting. Let's think about possible scenarios.} \\
\\ 
\texttt{If I take the lamp (3 points), then the remaining items are 3 thimbles and 1 rope. The other player could take all 3 thimbles and the rope. If their value for thimbles is low, maybe they'd prefer the rope. But if they value thimbles more, taking all three might be good. However, since I don't know their values, I need to make a proposal that's flexible.} \\
\\ 
\texttt{Alternatively, if I take two thimbles (2*2=4) and the rope (1), that's 5 points. Then the other player could take the lamp and one thimble. But if they value the lamp highly, that might work. But if they also want thimbles, maybe they'd prefer more of them. However, there are only three, so if I take two, they can take one. But maybe they want the lamp and the rope instead.} \\
\\ 
\texttt{Another option is to take the lamp and one thimble (3+2=5). Then the other player could take two thimbles and the rope. But again, depends on their values.} \\
\\ 
\texttt{Wait, the key is to make a proposal that doesn't conflict with possible other proposals. Since the other player's values are unknown, I need to choose a proposal that allows them to have a set that doesn't overlap with mine. The safest way might be to take the lamp and leave the thimbles and rope for them. Because if they want the lamp, we'd conflict, but if they don't, they can take other items. Alternatively, if I take the rope, which I value low, maybe that leaves the lamp for them if they value it.} \\
\\ 
\texttt{But my highest value is the lamp. So maybe I should go} \\
            }
        }
    }
    & & \\ \\

    \theutterance \stepcounter{utterance}  
    & & \multicolumn{4}{p{0.6\linewidth}}{
        \cellcolor[rgb]{0.9,0.9,0.9}{
            \makecell[{{p{\linewidth}}}]{
                \texttt{\tiny{[P1$\langle$GM]}}
                \texttt{<think>} \\
\texttt{Okay, let's try to figure out the best way to approach this. The available items are 1 lamp, 3 thimbles, and 1 rope. My values are lamp: 3, thimble: 2, rope: 1. The other player's values are secret, so I need to make a proposal that maximizes my points while allowing them to have a complementary set.} \\
\\ 
\texttt{First, I should consider what the other player might value. Since there's only one lamp and one rope, they might be high{-}value items for them too. But with three thimbles, maybe they can take some without conflicting. Let's think about possible scenarios.} \\
\\ 
\texttt{If I take the lamp (3 points), then the remaining items are 3 thimbles and 1 rope. The other player could take all 3 thimbles and the rope. If their value for thimbles is low, maybe they'd prefer the rope. But if they value thimbles more, taking all three might be good. However, since I don't know their values, I need to make a proposal that's flexible.} \\
\\ 
\texttt{Alternatively, if I take two thimbles (2*2=4) and the rope (1), that's 5 points. Then the other player could take the lamp and one thimble. But if they value the lamp highly, that might work. But if they also want thimbles, maybe they'd prefer more of them. However, there are only three, so if I take two, they can take one. But maybe they want the lamp and the rope instead.} \\
\\ 
\texttt{Another option is to take the lamp and one thimble (3+2=5). Then the other player could take two thimbles and the rope. But again, depends on their values.} \\
\\ 
\texttt{Wait, the key is to make a proposal that doesn't conflict with possible other proposals. Since the other player's values are unknown, I need to choose a proposal that allows them to have a set that doesn't overlap with mine. The safest way might be to take the lamp and leave the thimbles and rope for them. Because if they want the lamp, we'd conflict, but if they don't, they can take other items. Alternatively, if I take the rope, which I value low, maybe that leaves the lamp for them if they value it.} \\
\\ 
\texttt{But my highest value is the lamp. So maybe I should go} \\
\\ 
\\ 
\texttt{The time is up. It is now your turn to submit a secret proposal.} \\
            }
        }
    }
    & \\ \\

    \theutterance \stepcounter{utterance}  
    & \multicolumn{4}{p{0.6\linewidth}}{
        \cellcolor[rgb]{0.9,0.9,0.9}{
            \makecell[{{p{\linewidth}}}]{
                \texttt{\tiny{[P1$\rangle$GM]}}
                \texttt{<think>} \\
\texttt{Okay, let's try to figure out the best way to approach this. The available items are 1 lamp, 3 thimbles, and 1 rope. My values are lamp: 3, thimble: 2, rope: 1. The other player's values are secret, so I need to make a proposal that maximizes my points while allowing them to have a complementary set.} \\
\\ 
\texttt{First, I should consider what the other player might value. Since there's only one lamp and one rope, they might be high{-}value items for them too. But with three thimbles, maybe they can take some without conflicting. Let's think about possible scenarios.} \\
\\ 
\texttt{If I take the lamp (3 points), then the remaining items are 3 thimbles and 1 rope. The other player could take all 3 thimbles and the rope. If their value for thimbles is low, maybe they'd prefer the rope. But if they value thimbles more, taking all three might be good. However, since I don't know their values, I need to make a proposal that's flexible.} \\
\\ 
\texttt{Alternatively, if I take two thimbles (2*2=4) and the rope (1), that's 5 points. Then the other player could take the lamp and one thimble. But if they value the lamp highly, that might work. But if they also want thimbles, maybe they'd prefer more of them. However, there are only three, so if I take two, they can take one. But maybe they want the lamp and the rope instead.} \\
\\ 
\texttt{Another option is to take the lamp and one thimble (3+2=5). Then the other player could take two thimbles and the rope. But again, depends on their values.} \\
\\ 
\texttt{Wait, the key is to make a proposal that doesn't conflict with possible other proposals. Since the other player's values are unknown, I need to choose a proposal that allows them to have a set that doesn't overlap with mine. The safest way might be to take the lamp and leave the thimbles and rope for them. Because if they want the lamp, we'd conflict, but if they don't, they can take other items. Alternatively, if I take the rope, which I value low, maybe that leaves the lamp for them if they value it.} \\
\\ 
\texttt{But my highest value is the lamp. So maybe I should go} \\
            }
        }
    }
    & & \\ \\

    \theutterance \stepcounter{utterance}  
    & & & \multicolumn{2}{p{0.3\linewidth}}{
        \cellcolor[rgb]{0.9,0.9,0.9}{
            \makecell[{{p{\linewidth}}}]{
                \texttt{\tiny{[GM$|$GM]}}
                \texttt{continue} \\
            }
        }
    }
    & & \\ \\

    \theutterance \stepcounter{utterance}  
    & & & \multicolumn{2}{p{0.3\linewidth}}{
        \cellcolor[rgb]{0.9,0.9,0.9}{
            \makecell[{{p{\linewidth}}}]{
                \texttt{\tiny{[GM$|$GM]}}
                \texttt{player was instructed to make a proposal} \\
            }
        }
    }
    & & \\ \\

\end{supertabular}
}

\end{document}
